% Generated by scripts/gen_blueprint.py from blueprint/nodes/07-waldschmidt-1973.toml. Do not edit.

\chapter{Waldschmidt 1973 in full, and Schneider's eighth problem}
\label{chap:07-waldschmidt-1973}

Waldschmidt's Théorème of 1973 (Theorem~\ref{thm:two_algebraically_independent_of_exp_column}): if one column $e^{x_1y_2}$, $e^{x_2y_2}$ of the four exponentials is algebraic, then two of the eight numbers $x_i$, $y_j$, $e^{x_iy_j}$ are algebraically independent. Brownawell found it independently. Its proof reuses 33 results of the four exponentials development unchanged. The lemmas below carry the column case, where $e^{x_1y_1}$ and $e^{x_2y_1}$ are only algebraic over $\Q(\omega)$, so that their powers cost degree in $\omega$; they are sampled along the short side of the interpolation grid, where the degree budget absorbs them. Schneider's eighth problem, which gave the 1973 paper its title, follows (Theorem~\ref{thm:exp_e_or_exp_e_sq_transcendental}).

The notation is that of Chapter~\ref{chap:06-four-exponentials}. The \emph{column hypotheses} on $x_1,x_2,y_1,y_2\in\C$ are: $x_1,x_2$ are linearly independent over $\Q$, and so are $y_1,y_2$; $e^{x_1y_2}$ and $e^{x_2y_2}$ are algebraic; and
\[ \trdeg_{\Q}\Q[x_1,x_2,y_1,y_2,e^{x_1y_1},e^{x_2y_1}]\le1. \]

\section{One algebraic column}

% Prove2Me: FourExp.trdeg_one_presentation_column -- https://prove2.me/theorems/97fdb3d7-a247-4be9-a6ac-2b2b9b44b0cb
\begin{lemma}[Presenting the field, one algebraic column]
\label{lem:trdeg_one_presentation_column}
\lean{Diaz.trdeg_one_presentation_column}
\leanok
If $x_1,x_2,y_1,y_2$ satisfy the column hypotheses, then they have presentation data.

\emph{Source:} The reduction to $\omega$ in the proof of the Théorème, \cite[p.~196]{W73}.
\end{lemma}

\begin{proof}
\leanok
\uses{thm:hermite_lindemann_holds,lem:exists_monic_integral_model_presentation}
Take $\omega=x_1y_2$, which is transcendental by Theorem~\ref{thm:hermite_lindemann_holds}. All eight numbers are algebraic over $\Q(\omega)$, and Lemma~\ref{lem:exists_monic_integral_model_presentation} presents them at once.
\end{proof}

% Prove2Me: FourExp.exists_iteratedDeriv_presentation_of_exp_factor -- https://prove2.me/theorems/75faa454-0195-48c9-8276-574ade5ce1b0
\begin{lemma}[Presentations from the exponential factor]
\label{lem:exists_iteratedDeriv_presentation_of_exp_factor}
\lean{Diaz.exists_iteratedDeriv_presentation_of_exp_factor}
\leanok
Let $x_1,x_2,y_1,y_2,\omega,\omega_1\in\C$, let $\varphi$ be evaluation at $(\omega,\omega_1)$, and let $D,E_i,G_j\in\Z[X][Y]$ with $x_i\varphi(D)=\varphi(E_i)$ and $y_j\varphi(D)=\varphi(G_j)$. There is $c\in\N$ with the following property. Let $S,T,a,b,m,L,\delta\in\N$ and $\Lambda_1\in\C$, and let $U_{jh}\in\Z[X][Y]$ ($j,h<T$), of length at most $L$ and of degree at most $\delta$ in $X$ and in $Y$, satisfy
\[ \varphi(U_{jh})=\Lambda_1\,e^{(jx_1+hx_2)(ay_1+by_2)}. \]
Then there are $P_{ijh}\in\Z[X][Y]$ ($i<S$, $j,h<T$), of length at most $c^{c(1+m+S)}(1+m+S+T+a+b)^{c(1+m+S)}L$ and of degree at most $c(1+m+S)+\delta$ in $X$ and in $Y$, such that for all $f_{ijh}\in\C$
\[ \varphi(D)^{m+S}\,\Lambda_1\cdot\frac{d^{m}}{dz^{m}}\Bigl(\sum_{i,j,h}f_{ijh}\,z^{i}e^{(jx_1+hx_2)z}\Bigr)\Big|_{z=ay_1+by_2}=\sum_{i,j,h}f_{ijh}\,\varphi(P_{ijh}). \]

\emph{Source:} The presentation step that the proofs of \cite[Lemmes 4 and 7]{W73} share between the four exponentials case and the case of one algebraic column; the derivatives are expanded as in \cite[(7), p.~197]{W73}.
\end{lemma}

\begin{proof}
\leanok
\uses{lem:length_mul_le,lem:length_sum_le}
Leibniz's rule; $\varphi(D)^{m+S}$ clears the denominators of $jx_1+hx_2$ and of the powers of $ay_1+by_2$. Lengths are tracked with Lemmas~\ref{lem:length_sum_le} and~\ref{lem:length_mul_le}.
\end{proof}

% Prove2Me: FourExp.exists_iteratedDeriv_presentation_column -- https://prove2.me/theorems/2c076a9f-3bad-4cb7-ad6d-a247fba511b3
\begin{lemma}[Derivatives at a grid point, one algebraic column]
\label{lem:exists_iteratedDeriv_presentation_column}
\lean{Diaz.exists_iteratedDeriv_presentation_column}
\leanok
Let $x_1,x_2,y_1,y_2\in\C$ with $e^{x_1y_2}$ and $e^{x_2y_2}$ algebraic, and let $\omega,\omega_1$, $\varphi$ and $D,E_i,G_j,H_{ij}$ be as in Lemma~\ref{lem:exists_iteratedDeriv_presentation}. There is $c\in\N$ such that for all $S,T,a,b,m\in\N$ there are $\Lambda\neq0$ with $|\Lambda|\le c^{\,c(1+m+S+T(a+b))}$ and $P_{ijh}\in\Z[X][Y]$ ($i<S$, $j,h<T$), of length at most $c^{\,c(1+m+S+T(a+b))}(1+m+S+T+a+b)^{c(1+m+S)}$ and of degree at most $c(1+m+S+Ta)$ in $X$ and in $Y$, for which the identity of Lemma~\ref{lem:exists_iteratedDeriv_presentation} holds for all $f_{ijh}\in\C$.

\emph{Source:} A step of the proofs of \cite[Lemmes 4 and 7]{W73}, pp.~197 and 200.
\end{lemma}

\begin{proof}
\leanok
\uses{lem:exists_iteratedDeriv_presentation_of_exp_factor,lem:exists_int_pow_repr,lem:length_mul_le,lem:length_sum_le}
At $ay_1+by_2$ the exponential factor is $(e^{x_1y_1})^{ja}(e^{x_1y_2})^{jb}(e^{x_2y_1})^{ha}(e^{x_2y_2})^{hb}$. The powers of the algebraic column reduce through integer relations (Lemma~\ref{lem:exists_int_pow_repr}) and cost only height; the powers of $e^{x_1y_1}$ and $e^{x_2y_1}$, written through $H_{11}$ and $H_{21}$, cost degree proportional to $Ta$. Lemma~\ref{lem:exists_iteratedDeriv_presentation_of_exp_factor} does the rest.
\end{proof}

% Prove2Me: FourExp.exists_iteratedDeriv_reduced_presentation_column -- https://prove2.me/theorems/57c79d91-3641-4ba6-9b89-96a35793766f
\begin{lemma}[Reduced presentations, one algebraic column]
\label{lem:exists_iteratedDeriv_reduced_presentation_column}
\lean{Diaz.exists_iteratedDeriv_reduced_presentation_column}
\leanok
In the setting of Lemma~\ref{lem:exists_iteratedDeriv_presentation_column}, let moreover $Q\in\Z[X][Y]$ be monic in $Y$ with $\varphi(Q)=0$. Then the conclusion of Lemma~\ref{lem:exists_iteratedDeriv_reduced_presentation} holds with the $X$-degree bound $M+c(1+m+S+Ta)$ in place of $M+c(1+m+S)$.

\emph{Source:} A step of the proofs of \cite[Lemmes 4 and 7]{W73}, pp.~197 and 200.
\end{lemma}

\begin{proof}
\leanok
\uses{lem:exists_iteratedDeriv_presentation_column,lem:exists_modByMonic_length_le,lem:length_mul_le}
From Lemma~\ref{lem:exists_iteratedDeriv_presentation_column}, reducing modulo $Q$ with Lemma~\ref{lem:exists_modByMonic_length_le}.
\end{proof}

% Prove2Me: FourExp.aux_linear_system_column -- https://prove2.me/theorems/81adeeff-9cb0-4c16-90d4-33965d5c8b45
\begin{lemma}[The linear system, one algebraic column]
\label{lem:aux_linear_system_column}
\lean{Diaz.aux_linear_system_column}
\leanok
Let $x_1,x_2,y_1,y_2,\omega,\omega_1\in\C$ and $Q\in\Z[X][Y]$, and suppose that the derivatives of the functions $F_q$ have reduced presentations as in the conclusion of Lemma~\ref{lem:exists_iteratedDeriv_reduced_presentation_column}. Then the conclusion of Lemma~\ref{lem:aux_linear_system} holds, with $d=\deg_YQ$.

\emph{Source:} The equation count of \cite[Lemme 4]{W73}, pp.~197--198.
\end{lemma}

\begin{proof}
\leanok
\uses{lem:exists_pow_mul_pow_le_exp_sq_mul_sqrt_log}
As for Lemma~\ref{lem:aux_linear_system}, with the bounds of Lemma~\ref{lem:exists_pow_mul_pow_le_exp_sq_mul_sqrt_log}. Since $a<t_1$ gives $Ta\le2Nt_1\le2S$, a larger multiple of $S$ for $M$ still leaves twice as many unknowns as equations.
\end{proof}

% Prove2Me: FourExp.auxiliary_function_alg_column -- https://prove2.me/theorems/1d0ddadc-b140-4ea9-8c4f-b0cb29332ceb
\begin{lemma}[The auxiliary function, one algebraic column]
\label{lem:auxiliary_function_alg_column}
\lean{Diaz.auxiliary_function_alg_column}
\leanok
Let $x_1,x_2,y_1,y_2,\omega,\omega_1\in\C$ and $Q\in\Z[X][Y]$ with $d=\deg_YQ\ge1$, such that no non-zero $A\in\Z[X][Y]$ with $\deg_YA<d$ vanishes at $(\omega,\omega_1)$, and suppose that the derivatives of the functions $F_q$ have reduced presentations as in the conclusion of Lemma~\ref{lem:exists_iteratedDeriv_reduced_presentation_column}. Then the conclusion of Lemma~\ref{lem:auxiliary_function_alg} holds.

\emph{Source:} \cite[Lemme 4]{W73}, pp.~197--198, with Siegel's lemma, Lemme 1 there.
\end{lemma}

\begin{proof}
\leanok
\uses{lem:aux_linear_system_column,lem:siegel_aux}
Lemma~\ref{lem:aux_linear_system_column}, then Lemma~\ref{lem:siegel_aux}.
\end{proof}

% Prove2Me: FourExp.norm_to_polynomial_alg_column -- https://prove2.me/theorems/731b223a-6992-4c87-a410-1129b4aad5f9
\begin{lemma}[A small polynomial, one algebraic column]
\label{lem:norm_to_polynomial_alg_column}
\lean{Diaz.norm_to_polynomial_alg_column}
\leanok
Let $x_1,x_2,y_1,y_2,\omega,\omega_1\in\C$ and let $Q\in\Z[X][Y]$ be monic in $Y$ with $Q(\omega,\omega_1)=0$ and minimal there, and suppose that the derivatives of the functions $F_q$ have reduced presentations as in the conclusion of Lemma~\ref{lem:exists_iteratedDeriv_reduced_presentation_column}. Then the conclusion of Lemma~\ref{lem:norm_to_polynomial_alg} holds.

\emph{Source:} \cite[Lemme 7]{W73}, pp.~200--201.
\end{lemma}

\begin{proof}
\leanok
\uses{lem:exists_int_norm,lem:exists_pow_mul_pow_le_exp_sq_mul_sqrt_log,lem:length_sum_le,lem:length_mul_le}
As for Lemma~\ref{lem:norm_to_polynomial_alg}, with the norm of Lemma~\ref{lem:exists_int_norm}. Since $a<R_1$ gives $2Na\le28N^{2}/\sqrt{\log N}$, the $X$-degree stays $O(N^{2}/\sqrt{\log N})$ and only the constant $k$ grows.
\end{proof}

% Prove2Me: FourExp.construction_core_column -- https://prove2.me/theorems/690b22eb-5683-42a9-b0ee-fd8096209490
\begin{lemma}[The core of the construction, one algebraic column]
\label{lem:construction_core_column}
\lean{Diaz.construction_core_column}
\leanok
If $x_1,x_2,y_1,y_2$ satisfy the column hypotheses, then the conclusion of Lemma~\ref{lem:construction_core_1973} holds.

\emph{Source:} \cite[Lemmes 4, 5 and 7]{W73}, with the parameters of p.~196.
\end{lemma}

\begin{proof}
\leanok
\uses{lem:trdeg_one_presentation_column,lem:exists_iteratedDeriv_reduced_presentation_column,lem:auxiliary_function_alg_column,lem:extrapolation,lem:norm_to_polynomial_alg_column}
Lemmas~\ref{lem:trdeg_one_presentation_column}, \ref{lem:exists_iteratedDeriv_reduced_presentation_column}, \ref{lem:auxiliary_function_alg_column}, \ref{lem:extrapolation} and~\ref{lem:norm_to_polynomial_alg_column}, in this order. The numbers $e^{x_1y_1}$ and $e^{x_2y_1}$ are read only at the short index $a$ of the grid.
\end{proof}

% Prove2Me: FourExp.auxiliary_construction_column -- https://prove2.me/theorems/746d9340-01c6-4aa8-b08a-1c020dfaed36
\begin{lemma}[The construction, one algebraic column]
\label{lem:auxiliary_construction_column}
\lean{Diaz.auxiliary_construction_column}
\leanok
Let $x_1,x_2,y_1,y_2$ satisfy the column hypotheses. Then there are a transcendental $\omega$, and functions $\sigma_1,\sigma_2$ and constants $a_1,a_2\ge1$ satisfying the hypotheses of Theorem~\ref{thm:transcendence_criterion}, such that for every $C$ and every large $N$ there are natural numbers $S,T,R_1,R_2,S'$, coefficients $c(i,j,h)\in\C$ ($i<S$, $j,h<T$), not all zero, and $\lambda>0$ with the two properties listed in Lemma~\ref{lem:auxiliary_construction}, for these $x_1,x_2,y_1,y_2$.

\emph{Source:} \cite[Lemmes 4, 5 and 7]{W73}, with the growth functions of p.~201, and the count that \cite[Lemme 6]{W73} needs for the zero estimate \cite[\S4, Lemme 3]{W71}.
\end{lemma}

\begin{proof}
\leanok
\uses{lem:construction_growth,lem:construction_count_1973,lem:construction_core_column}
Lemma~\ref{lem:construction_core_column}, with the growth functions of Lemma~\ref{lem:construction_growth} and the count of Lemma~\ref{lem:construction_count_1973}.
\end{proof}

% Prove2Me: FourExp.small_polynomials_of_column -- https://prove2.me/theorems/4b4429d8-3091-408a-8b30-d661d6c0a9ef
\begin{lemma}[One algebraic column gives small polynomials]
\label{lem:small_polynomials_of_column}
\lean{Diaz.small_polynomials_of_column}
\leanok
Let $x_1,x_2,y_1,y_2$ satisfy the column hypotheses. Then there are a transcendental $\omega\in\C$, functions $\sigma_1,\sigma_2$ and constants $a_1,a_2\ge1$ satisfying the hypotheses of Theorem~\ref{thm:transcendence_criterion}, such that for every $C\in\R$ there is $N_0$ and for every $N>N_0$ a polynomial $P_N\in\Z[X]$ that is small at $\omega$ for $(N,C)$ with respect to $\sigma_1,\sigma_2$.

\emph{Source:} \cite[Lemmes 4--7]{W73}, with the zero estimate \cite[\S4, Lemme 3]{W71} in place of Gel'fond's zero lemma in Lemme 6.
\end{lemma}

\begin{proof}
\leanok
\uses{lem:auxiliary_construction_column,lem:nonvanishing_derivative}
Lemma~\ref{lem:auxiliary_construction_column}, with a non-zero derivative supplied by Lemma~\ref{lem:nonvanishing_derivative}.
\end{proof}

\section{The theorems}

% Prove2Me: DiazModulus.two_algebraically_independent_of_exp_column -- https://prove2.me/theorems/b3e7e62f-38e7-469f-9b86-446206bec1a3
\begin{theorem}[Waldschmidt 1973]
\label{thm:two_algebraically_independent_of_exp_column}
\lean{Diaz.two_algebraically_independent_of_exp_column}
\leanok
Let $x_1,x_2$ be complex numbers linearly independent over $\Q$, and $y_1,y_2$ complex numbers linearly independent over $\Q$. If $e^{x_1y_2}$ and $e^{x_2y_2}$ are algebraic, then two of the eight numbers
\[ x_1,\ x_2,\ y_1,\ y_2,\ e^{x_1y_1},\ e^{x_1y_2},\ e^{x_2y_1},\ e^{x_2y_2} \]
are algebraically independent over $\Q$.

\emph{Source:} Waldschmidt \cite[Théorème, p.~192]{W73}; found independently by Brownawell \cite{Bro74}.
\end{theorem}

\begin{proof}
\leanok
\uses{thm:hermite_lindemann_holds,lem:isAlgebraic_adjoin_of_not_algebraicIndependent_pair,lem:trdeg_adjoin_le_one_of_isAlgebraic_adjoin,lem:small_polynomials_of_column,thm:transcendence_criterion}
Suppose not. One of $x_1$, $y_2$ is transcendental by Theorem~\ref{thm:hermite_lindemann_holds}, since $x_1y_2\neq0$ and $e^{x_1y_2}$ is algebraic. Lemma~\ref{lem:isAlgebraic_adjoin_of_not_algebraicIndependent_pair} makes all eight numbers algebraic over the ring it generates, so by Lemma~\ref{lem:trdeg_adjoin_le_one_of_isAlgebraic_adjoin} the column hypotheses hold. Lemma~\ref{lem:small_polynomials_of_column} then gives polynomials too small at a transcendental number, which Theorem~\ref{thm:transcendence_criterion} makes algebraic.
\end{proof}

% Prove2Me: Transcendence.exp_e_or_exp_e_sq_transcendental -- https://prove2.me/theorems/88f16d63-3fc8-487c-a4fe-f9b72badba3e
\begin{theorem}[Schneider's eighth problem]
\label{thm:exp_e_or_exp_e_sq_transcendental}
\lean{Transcendence.exp_e_or_exp_e_sq_transcendental}
\leanok
At least one of the numbers $e^{e}$ and $e^{e^{2}}$ is transcendental.

\emph{Source:} The problem: Schneider \cite{Sch57}, as quoted in \cite[p.~191]{W73}. The solution: \cite[p.~192]{W73}, the case $r=1$ of the solution of Schneider's problem, deduced from Corollaire 1; found independently by Brownawell \cite{Bro74}.
\end{theorem}

\begin{proof}
\leanok
\uses{thm:hermite_lindemann_holds,lem:trdeg_adjoin_le_one_of_isAlgebraic_adjoin,thm:two_algebraically_independent_of_exp_column}
The number $e=e^{1}$ is transcendental by Theorem~\ref{thm:hermite_lindemann_holds}. Apply Theorem~\ref{thm:two_algebraically_independent_of_exp_column} at $x=y=(1,e)$: the column $y_2=e$ carries $e^{e}$ and $e^{e^{2}}$, and if both were algebraic, all eight numbers would be algebraic over $\Q[e]$, of transcendence degree one by Lemma~\ref{lem:trdeg_adjoin_le_one_of_isAlgebraic_adjoin}.
\end{proof}
